\scanpage{189}
\mspage{63}
% archive sheet number: 188

\deleted{part, si $\deleted{\unsure{X}}Z\longrightarrow Y,Z'\unsure{-|\longrightarrow}Y$ sont dans $M^{\unsure{\pi|\Pi}}/Y$, et si on a une double flèche}

\deleted{\diagram{\begin{tikzcd}[ampersand replacement=\&,column sep=small,row sep=small]
 \deleted{Z}\arrow[rr,shift left=.6ex]\arrow[rr,shift right=.6ex]\arrow[dr] \& \& \deleted{Z'}\arrow[dl]\\
 \& \deleted{Y} \&
\end{tikzcd}}}

\deleted{en factorisant par $Z'\longrightarrow\underset{\substack{\rotatebox{-90}{$\in$}\\\deleted{\unsure{M}}\mathrm{Ss}^{\pi}_{\pi}(Y)\subset M^\pi/Y}}{Z_\alpha}\lhook\joinrel\longrightarrow Y$, on trouve $Z'\longrightarrow Z_\alpha$ qui égalise la double flèche.}

Considérons le foncteur canonique

\[
 I=\mathrm{Ss}_{\pi}(Y)\xrightarrow{\unsure{\varphi|\phi}}\underline M^\pi/Y=\underset{\unsure{9|g|\varphi}}{I'},
\]

où la catégorie source \inserted{$I$} est filtrante. Par SGA 4 I 8.1.3, ce foncteur est cofinal sss il satisfait deux conditions \unsure{F1, F2|F1, F7}, et alors \deleted{\unsure{$I$}} $I'$ est également filtrante. La \deleted{\unsure{même}} condition F1, \deleted{\unsure{dit que}} est la condition habituelle que tt $i'\in\mathrm{Ob}\,I'$ est ``majoré'' par un $\unsure{\varphi(i)}$, i.e. \unsure{$\forall$} $Z\longrightarrow Y$, avec $Z\in\underline M^\pi$, est ``majoré'' par un \inserted{\uline{mono}} $Z_\alpha\hookrightarrow Y$, \inserted{avec $Z_\alpha\in M^\pi$} ce qui résulte du fait que \deleted{si $Z$} c'est majoré par $Z_\alpha\lhook\joinrel\longrightarrow Y$ avec $Z_\alpha\in\mathrm{Filt}^{\pi}$, \deleted{et que} i.e $Z_\alpha\in\mathrm{Ss}_\pi(Y)$, et que parmi ces $Z_\alpha$ ceux qui sont dans $M^\pi$ sont cofinaux lorsque $Y\in M$.

% slanted
\margin{Comme $I$ est filtrante}

L'autre condition \unsure{*} est qu'une double flèche $i'\mathrel{\substack{\xrightarrow{f',g'}\\[-.5ex]\longrightarrow}}\varphi(i)$ \unsure{puisse|pour}

